% ======================================================================
%  Section 4.3.3 --- Limitation: the estimator is structure-blind.
%  Corrected.
%
%  The argument is sound and the permutation-invariance point is the best
%  single sentence in the section. Four changes.
%
%  A. "the multi-scale Hermite responses ALREADY COMPUTED by the scheme
%     provide direct access to it" -- the same claim corrected in the
%     introduction and flagged for 4.4.1. After the move to first-order
%     kernels in 4.4.2 the structure measure is NOT among the responses the
%     feature vector uses; the implementation evaluates six additional
%     separable convolutions per iteration. This is the third place the
%     claim appears and all three must agree.
%
%  B. The central prediction of the second paragraph is now measured.
%     Section 4.7.3 reports that at iteration 1, on a CLEAN scene where
%     nothing is corrupted, 13.6% of measurements carry a weight below a
%     quarter -- every one of them valid. That is this paragraph's
%     prediction, quantified. It should say so.
%
%  C. The closing sentence promises that the structural information
%     "resolves the ambiguity". The results do not establish that. Table 4.3
%     reports 13.6% early over-rejection for VARIANT 3, with the structure
%     gate active, so the limitation identified here survives the proposed
%     remedy in the only configuration measured. Promise the attempt, not
%     the outcome.
%
%  D. LABEL. My correction to 4.3.1 refers to this subsection as
%     \ref{sec:structure-blind}. Your label is sec:tukey-limitation. Keep
%     yours -- it is the one already in the document -- and change the
%     reference in 4.3.1 to match. Do not add a second label.
% ======================================================================

\subsection{Limitation: the estimator is structure-blind}
\label{sec:tukey-limitation}

The weighting~\eqref{eq:tukey-weights} judges each measurement exclusively by
the magnitude of its own residual. It has no access to the spatial arrangement
of the measurements, and would assign identical weights to a residual vector
and to any permutation of it, even though the measurements corrupted by an
occlusion form a spatially coherent block rather than an arbitrary scattered
subset. Nor does it have access to the local image structure at the position
of each measurement, and therefore cannot distinguish a large residual arising
where the image is photometrically sensitive to camera motion, and which is
correspondingly informative about the pose error, from an equally large
residual arising where the local structure contributes little to the
interaction matrix.

The consequence is a systematic difficulty at the beginning of the servo. When
the initial pose error is large, valid measurements in photometrically
sensitive regions produce legitimately large residuals, which a redescending
estimator may reject; conversely, measurements whose residual happens to be
moderate are retained irrespective of how much they constrain the motion.
%%% (B) The prediction, measured. This is the strongest support the section
%%% can have and it costs two sentences. Note that it is measured on a CLEAN
%%% scene, which is what makes it unambiguous: nothing is corrupted, so
%%% every rejection is an error.
The effect is not hypothetical. Section~\ref{sec:weight-maps} measures it
directly on an unoccluded target, where by construction no measurement is
corrupted and every rejection is therefore mistaken: at the first iteration
$13.6\%$ of measurements carry a weight below a quarter and $3.5\%$ are
rejected outright. The estimator discards roughly one measurement in seven
purely because the pose error is still large, and recovers them only as the
residual distribution contracts.

The information required to resolve this ambiguity is absent from the
residuals but present in the local image structure, which the Hermite
representation the scheme already operates in makes available at the analysis
scales the controller uses. The following section uses it to define a
structure-aware confidence measure; whether doing so mitigates the difficulty
described here is an empirical question, answered in
Section~\ref{sec:weight-maps}.
%%% (A) "already computed by the scheme" removed: the measure is obtained
%%% from the first-order kernels of the same family at the same scales, at
%%% the cost of six additional convolutions per iteration. The wording above
%%% claims only that the REPRESENTATION is already in use, which is true.
%%%
%%% (C) "resolve this ambiguity" softened to "mitigates ... is an empirical
%%% question". Table 4.3's 13.6% is Variant 3's own weight map, so the
%%% limitation identified here is still present when the proposed remedy is
%%% active. Claiming resolution here would be contradicted three sections
%%% later by the chapter's own table.

% ======================================================================
%  E. THE MEASUREMENT THAT WOULD SETTLE THIS -- worth an hour
%
%  Table 4.3 shows 13.6% early over-rejection for Variant 3. What it cannot
%  show is whether that is better or worse than Tukey alone, because no
%  weight map was produced for Variant 2. Without it the chapter cannot say
%  whether the structure gate mitigates the difficulty of this section at
%  all -- and that is the chapter's central mechanism claim, which the pose
%  errors of 4.7.2 left as a tie.
%
%  The measurement is cheap: add saveMaps() to Variant 2 and run nominal
%  once. Two numbers decide it, the fraction below a quarter weight at
%  iteration 1 for each variant:
%
%     V3 (measured):  13.6%
%     V2 (unknown):      ?%
%
%  If V2 is substantially higher, the structure gate demonstrably reduces
%  the over-rejection this section predicts, and that is a result the
%  chapter currently lacks. If the two are the same, say so -- it is the
%  honest counterpart of the tie in 4.7.2 and it tells the reader where the
%  mechanism does and does not act.
%
%  This is the single highest-value run still outstanding in the chapter.
% ======================================================================
